\section{Theory and Methods}
\label{sec:theory}

% ======================================================================
\subsection{XC Functional Construction}
\label{sec:theory-dft}

For a molecular system with density $\rho$, the Kohn--Sham energy is
written as
\begin{equation}
  E_\theta[\rho]
  =
  T_s[\rho]
  + E_{\rm ne}[\rho]
  + J[\rho]
  + E_{\rm xc}^{\theta}[\rho]
  + E_{\rm nn},
  \label{eq:dft-energy}
\end{equation}
where $T_s$ is the noninteracting kinetic energy, $E_{\rm ne}$ is the
electron--nuclear attraction, $J$ is the Hartree energy, $E_{\rm nn}$
is the nuclear repulsion, and $\theta$ denotes the neural XC
parameters. In an AO basis, the corresponding generalized
eigenvalue problem is
\begin{align}
  \mathbf{F}_\theta[\mathbf{D}]\mathbf{C}
  &=
  \mathbf{S}\mathbf{C}\boldsymbol{\epsilon},
  \label{eq:ks-fock}
  \\
  \mathbf{F}_\theta
  &=
  \mathbf{h}
  + \mathbf{J}[\mathbf{D}]
  + \mathbf{V}_{\rm xc}^{\theta}[\mathbf{D}]
  + \mathbf{V}_{\rm hyb}^{\theta}[\mathbf{D}].
  \nonumber
\end{align}
where $\mathbf{D}$ is the density matrix and
$\mathbf{V}_{\rm hyb}^{\theta}$ collects nonlocal hybrid
contributions.

GradTDDFT uses a coefficient-basis neural hybrid functional. Following
the energy-density mixture idea used in GradDFT,\cite{GradDFT} the XC
energy is written in continuous form as
\begin{align}
  E_{\rm xc}^{\theta}
  &=
  \int e_{\rm xc}^{\theta}(\mathbf{r})\,d\mathbf{r},
  \nonumber
  \\
  e_{\rm xc}^{\theta}(\mathbf{r})
  &=
  \sum_{m=1}^{M}
  c_m^{\theta}(\mathbf{z}(\mathbf{r}))e_m^{\rm sl}(\mathbf{r})
  \nonumber
  \\
  &\quad
  + c_x^{\theta}(\mathbf{z}(\mathbf{r}))h(\mathbf{r})
  + c_2^{\theta}(\mathbf{z}(\mathbf{r}))p(\mathbf{r}) .
  \label{eq:neural-hybrid-functional}
\end{align}
The energy density is expanded into three channel groups: local or
semi-local XC energy densities, a projected local exact-exchange
channel, and an optional projected local MP2/PT2 second-order
correlation channel. Here $\mathbf{z}(\mathbf{r})$ is a local
descriptor vector, $c_m^\theta$, $c_x^\theta$, and $c_2^\theta$ are
learned pointwise mixing fields, $e_m^{\rm sl}(\mathbf{r})$ are
semi-local energy-density channels, $h(\mathbf{r})$ is the local HF
exchange channel density, and $p(\mathbf{r})$ is the local PT2 channel
density. Constant coefficients recover conventional hybrid or
double-hybrid forms, while learned pointwise coefficients define the
neural functional used in GradTDDFT. In numerical calculations,
Eq.~\eqref{eq:neural-hybrid-functional} is evaluated by molecular
quadrature.

A general semi-local channel has the form
\begin{equation}
  e_m^{\rm sl}(\mathbf{r})
  =
  \rho(\mathbf{r})\,
  \varepsilon_m^{\rm sl}
  \!\left(
    \rho_\sigma,
    \nabla\rho_\sigma\!\cdot\!\nabla\rho_{\sigma'},
    \tau_\sigma
  \right),
  \label{eq:semilocal-channel}
\end{equation}
with LDA, GGA, and meta-GGA channels obtained by selecting the
corresponding subset of density variables. For example, a
spin-unpolarized LDA exchange channel and a GGA exchange channel may
be written as
\begin{align}
  e_x^{\rm LDA}(\mathbf{r})
  &=
  -C_x\rho(\mathbf{r})^{4/3},
  \nonumber
  \\
  e_x^{\rm GGA}(\mathbf{r})
  &=
  e_x^{\rm LDA}(\mathbf{r})F_x(s),
  \nonumber
  \\
  s=
  &\,
  \frac{|\nabla\rho(\mathbf{r})|}
       {2(3\pi^2)^{1/3}\rho(\mathbf{r})^{4/3}},
  \label{eq:lda-gga-channel-example}
\end{align}
where $C_x=\frac{3}{4}(3/\pi)^{1/3}$ and $F_x$ is the
functional-specific GGA enhancement factor. Through jax-xc
\cite{JAXXC}, different conventional LDA, GGA, and meta-GGA
energy-density forms can be inserted into the semi-local channel set
without changing the SCF or TDDFT equations.

The HF and MP2/PT2 ingredients enter the coefficient basis through
local channel densities whose integrals recover the conventional
nonlocal energies.  With the spin-density matrix
$\gamma_\sigma(\mathbf{r},\mathbf{r}')$, the exact-exchange channel is
\begin{align}
  E_x^{\rm HF}
  &=
  \int h(\mathbf{r})\,d\mathbf{r}
  =
  -\frac{1}{2}
  \sum_\sigma
  \int
  \frac{
    |\gamma_\sigma(\mathbf{r},\mathbf{r}')|^2
  }{
    |\mathbf{r}-\mathbf{r}'|
  }
  \,d\mathbf{r}\,d\mathbf{r}'
  \nonumber
  \\
  &=
  -\frac{1}{2}
  \sum_{\sigma}
  \sum_{ij\in{\rm occ}_\sigma}
  (ij|ji).
  \label{eq:hf-energy-main}
\end{align}
For the second-order correlation channel, using spin-orbital indices
$i,j$ for occupied orbitals and $a,b$ for virtual orbitals,
\begin{align}
  E_c^{(2)}
  &=
  \int p(\mathbf{r})\,d\mathbf{r}
  =
  -\frac{1}{4}
  \sum_{ijab}
  \frac{
    |\langle ij||ab\rangle|^2
  }{
    D_{ij}^{ab}
  },
  \label{eq:pt2-energy-main}
  \\
  D_{ij}^{ab}
  &=
  \epsilon_a+\epsilon_b-\epsilon_i-\epsilon_j,
  \qquad
  \langle ij||ab\rangle=(ia|jb)-(ib|ja).
  \nonumber
\end{align}
Here $(pq|rs)$ denotes a two-electron integral in chemists'
notation.  The local gauges $h(\mathbf{r})$ and $p(\mathbf{r})$ are
not unique; in GradTDDFT they are projected to the molecular grid as
local HF and PT2 channel features.  Their detailed grid definitions
and response binding are given in the Supporting Information.

% ======================================================================
\subsection{Differentiable SCF and TDDFT}
\label{sec:tddft}

GradTDDFT treats the molecular integrals, basis metadata, and grid
quantities as fixed input data $\mathcal{I}$. The density-dependent
calculation is written as a differentiable SCF fixed point,
\begin{equation}
  \mathbf{D}_\theta^\ast
  =
  \Phi_\theta(\mathbf{D}_\theta^\ast;\mathcal{I}),
  \qquad
  \mathbf{R}_\theta(\mathbf{D})
  =
  \mathbf{D}-\Phi_\theta(\mathbf{D};\mathcal{I}),
  \label{eq:scf-fixed-point-main}
\end{equation}
where $\Phi_\theta$ denotes Fock construction, generalized
diagonalization, occupation assignment, and density reconstruction.
For implicit self-consistent differentiation, the converged density
satisfies $\mathbf{R}_\theta(\mathbf{D}_\theta^\ast)=0$, and its
parameter derivative is obtained from
\begin{equation}
  \frac{d\mathbf{D}_\theta^\ast}{d\theta}
  =
  -
  \left(
    \frac{\partial\mathbf{R}_\theta}{\partial\mathbf{D}}
  \right)^{-1}
  \frac{\partial\mathbf{R}_\theta}{\partial\theta}.
  \label{eq:implicit-main}
\end{equation}
This implicit SCF relation is the ground-state differential used by
the self-consistent training path; alternative fixed-density and
finite-SCF differentiation modes are defined in the Supporting
Information.

After SCF convergence, excitation energies are obtained from
linear-response TDDFT. For occupied orbitals $i,j$ and virtual
orbitals $a,b$, the full Casida equation is
\begin{equation}
  \begin{pmatrix}
    \mathbf{A} & \mathbf{B} \\
    \mathbf{B} & \mathbf{A}
  \end{pmatrix}
  \begin{pmatrix}
    \mathbf{X}_I \\
    \mathbf{Y}_I
  \end{pmatrix}
  =
  \Omega_I
  \begin{pmatrix}
    \mathbf{1} & \mathbf{0} \\
    \mathbf{0} & -\mathbf{1}
  \end{pmatrix}
  \begin{pmatrix}
    \mathbf{X}_I \\
    \mathbf{Y}_I
  \end{pmatrix},
  \label{eq:casida-main}
\end{equation}
where $\Omega_I$ is the excitation energy. The response matrices have
the schematic form
\begin{equation}
  \begin{aligned}
  A_{ia,jb}
  &=
  (\epsilon_a-\epsilon_i)\delta_{ij}\delta_{ab}
  + K_{ia,jb},
  \\
  B_{ia,jb}
  &=
  K_{ia,bj},
  \end{aligned}
  \label{eq:ab-main}
\end{equation}
where $K$ contains Coulomb, XC-kernel, and hybrid-response
contributions. The Tamm--Dancoff approximation sets
$\mathbf{B}=0$ and gives
\begin{equation}
  \mathbf{A}\mathbf{X}_I=\Omega_I\mathbf{X}_I .
  \label{eq:tda-main}
\end{equation}

The differentiable link between the SCF energy and the TDDFT response
is the XC kernel. For the local and semi-local part,
\begin{equation}
  f_{\rm xc}^{\theta}(\mathbf{r},\mathbf{r}')
  =
  \frac{\delta^2 E_{\rm xc}^{\theta}}
       {\delta\rho(\mathbf{r})\delta\rho(\mathbf{r}')}.
  \label{eq:fxc-main}
\end{equation}
Thus the SCF potential and the local response kernel are derivatives
of the same energy-density model. In GradTDDFT these derivatives are
evaluated by automatic differentiation of the pointwise energy
density, so derivatives of the learned mixing fields are included in
both the potential and the TDDFT kernel. The nonlocal HF response,
the PT2 excited-state correction, and optional RI response
contractions are specified in the Supporting Information.

For differentiable excited-state training, the response problem is
evaluated at the implicit SCF solution $\mathbf{D}_\theta^\ast$. The
TDDFT root derivative therefore contains both the explicit derivative
of the response operator and the implicit derivative of the converged
density in Eq.~\eqref{eq:implicit-main}. In the Davidson response
solver, GradTDDFT follows the implicit eigensolver differentiation
strategy of Xie, Liu, and Wang,\cite{Xie2020DominantEigensolverAD}
using the converged Ritz amplitudes to reconstruct the root-energy
differential rather than differentiating through the full iteration
history. For the TDA operator, this takes the Rayleigh-quotient form
\begin{equation}
  d\Omega_I
  =
  \frac{
    \mathbf{X}_I^{T}
    (d\mathbf{A})
    \mathbf{X}_I
  }{
    \mathbf{X}_I^{T}\mathbf{X}_I
  },
  \label{eq:tda-root-differential-main}
\end{equation}
with $d\mathbf{A}$ including both direct neural-kernel derivatives and
the SCF response from Eq.~\eqref{eq:implicit-main}. The corresponding
full Casida expression is given in the Supporting Information.
